% packets on the unrolled torsion line: \packetrow{x}{y}{sigma in degrees}{width cm}{height cm} draws the three
% periodic Gaussians at tau = 0, 120, 240 (-120) over [-180,180], gold fill, overlap darker, profiles fine, union heavy.
\newcommand\packetrow[5]{\begin{scope}[shift={(#1,#2)},x={#4cm/360},y={#5cm}]
  \draw[fine] (-180,0)--(180,0);
  \foreach \c in {0,120,-120} {\path[fill=ColRetained!55] plot[domain=-180:180,samples=121,variable=\t] (\t,{exp(-((\t-\c)/#3)^2/2)+exp(-((\t-\c-360)/#3)^2/2)+exp(-((\t-\c+360)/#3)^2/2)}) -- (180,0) -- (-180,0) -- cycle;}
  % overlap: the smaller of each neighbouring pair, drawn as the min of all pairs (two packets at most overlap visibly)
  \path[fill=ColRetained] plot[domain=-180:180,samples=181,variable=\t] (\t,{max(min(exp(-((\t)/#3)^2/2)+exp(-((\t-360)/#3)^2/2)+exp(-((\t+360)/#3)^2/2),exp(-((\t-120)/#3)^2/2)+exp(-((\t-480)/#3)^2/2)+exp(-((\t+240)/#3)^2/2)),min(exp(-((\t)/#3)^2/2)+exp(-((\t-360)/#3)^2/2)+exp(-((\t+360)/#3)^2/2),exp(-((\t+120)/#3)^2/2)+exp(-((\t-240)/#3)^2/2)+exp(-((\t+480)/#3)^2/2)),min(exp(-((\t-120)/#3)^2/2)+exp(-((\t-480)/#3)^2/2)+exp(-((\t+240)/#3)^2/2),exp(-((\t+120)/#3)^2/2)+exp(-((\t-240)/#3)^2/2)+exp(-((\t+480)/#3)^2/2)))}) -- (180,0) -- (-180,0) -- cycle;
  \foreach \c in {0,120,-120} {\draw[fine] plot[domain=-180:180,samples=121,variable=\t] (\t,{exp(-((\t-\c)/#3)^2/2)+exp(-((\t-\c-360)/#3)^2/2)+exp(-((\t-\c+360)/#3)^2/2)});}
  \foreach \c/\l in {0/{g_0},120/{g_1},-120/{g_2}} {\hatom{\c}{0}{.07}{ColDot}}
  \end{scope}}
